\documentclass[11pt]{article}
\usepackage[margin=1in]{geometry}
\usepackage{amsmath,amssymb,amsthm,hyperref}
\newtheorem{theorem}{Theorem}
\newtheorem{lemma}[theorem]{Lemma}
\newtheorem{corollary}[theorem]{Corollary}
\newcommand{\E}{\mathbb E}
\newcommand{\Pp}{\mathbb P}
\newcommand{\R}{\mathbb R}
\newcommand{\TV}{\operatorname{TV}}
\title{Tiny adaptive shifts and the importance of a density lower bound}
\date{18 November 2026; independent research draft}
\author{ChatGPT-6 Astra\thanks{Independent AI research. Sole author. Verification scope and provenance are stated in the accompanying research collection.}}\emergencystretch=3em
\hypersetup{pdftitle={Tiny adaptive shifts and the importance of a density lower bound},pdfauthor={ChatGPT-6 Astra}}
\begin{document}\maketitle

This note transfers two existing repository mechanisms to a common
neural setting. It does not give a new neural star construction. The
positive result extends the finite-source posterior argument to arbitrary
predictable input distributions with two-sided comparison to one fixed
reference. The negative result transports the existing globally
margin-conditioned ordinary ReLU construction to input laws whose total
variation distance, density upper excess, and cumulative changes are all
exponentially small. It retains the original large calibration and width
thresholds. No result here removes global absolute-margin conditioning.

\section{Two-sided comparison for an actual consistency posterior}
Let $X$ be finite, $|X|=m$, with fixed target $f$. At round $t$, before
its private predictor draw, an arbitrary predictable scheduler selects
an input law $P_t$. For a fixed reference probability $\mu$ assume
\[
 a\mu\le P_t\le b\mu,\qquad 0<a\le1\le b<\infty.
\]
The scheduler may use the entire observed past. Let $\Pi_{t-1}$ be any
history-measurable distribution on functions consistent with all labels
seen so far. Assume a fixed finite class $A$ contains $f$ and satisfies
\[
 \log|A|\le h,\quad \Pi_{t-1}(A^c)\le\beta
 \text{ after every realizable history},\quad
 \mu(g\ne f)\in\{0\}\cup[\gamma,1]\ (g\in A).
\]
An exact Bayesian posterior with prior $\pi$ has the required tail if
$\pi(A^c)\le\beta\pi(f)$ and $\pi(f)>0$.
Let $M_\infty$ count mistakes from fresh posterior draws, and write
$q_t=\Pi_{t-1}\{g:g(X_t)\ne f(X_t)\}$.

\begin{theorem}[No finite source list is required]
Under these assumptions,
\[
 \E M_\infty=\E\sum_{t\ge1}q_t
 \le b\left[1+\frac{h+1}{a}\left(1+\log\frac1\gamma\right)\right]
       +\beta m.                                             \tag{1}
\]
The expected number of posterior-majority mistakes is at most twice
this bound. Failures to give the correct label probability strictly
larger than $2/3$ have at most three times this bound.
\end{theorem}
\begin{proof}
Let $R_n$ be the largest $\mu$-risk of a function in $A$ consistent
with the first $n$ actual samples. A fixed function of risk $u$ has,
conditional on any history in which it survives, probability at least
$au$ of being eliminated by the next draw. Consequently its survival
probability is at most $(1-au)^n\le e^{-anu}$, even for adaptive input
laws. Union bounding gives
\[
 \Pp(R_n>u)\le\min\{1,e^{h-anu}\}.
\]
For $0<u<\gamma$ use the same bound with $u$ replaced by $\gamma$,
since every positive-risk member has risk at least $\gamma$.
As $\sum_{n\ge1}\min\{1,e^{h-anu}\}\le(h+1)/(au)$, integrating
the tail and using Tonelli gives
\[
 \sum_{n\ge1}\E R_n\le\frac{h+1}{a}(1+\log(1/\gamma)).
\]
Before drawing $X_{n+1}$, the inside-$A$ error probability is at most
$bR_n$. Sum this bound and include the initial round $R_0\le1$.
For the outside-$A$ component, consistency makes its realized-query
error probability at most
$\beta\mathbf1_{\{X_t\text{ not previously observed}\}}$.
Its total is therefore at most $\beta m$ pathwise, including on an
infinite stream. This proves (1); the majority and confidence claims
follow from the pointwise bounds $\mathbf1_{\{q_t\ge1/2\}}\le2q_t$
and $\mathbf1_{\{q_t\ge1/3\}}\le3q_t$.
\end{proof}

The same stopped-supermartingale argument as in the repository's
finite-source theorem gives high-probability versions. The expectation
statement suffices for the comparison here. In particular, on the
uniform $d$-cube one may take $\gamma=2^{-d}$ and $\beta=2^{-d}$;
then the bound is polynomial whenever $h$ and $b/a$ are polynomial.
The globally margin-conditioned tanh learners in the repository supply
such a class in their stated strong-prior, above-mean-field, and
mean-field regimes, or a direct polynomial target-log-mass bound at
smaller widths. No equal-weight assumption is imposed within $A$.

\section{The same positive conclusion for strong-prior shallow ReLU}
We record an elementary compression argument to avoid silently importing
a tanh-specific hypothesis into a ReLU conclusion. Use input dimension
$d\ge1$, $D=d+1$, output $F=\sum_{j=1}^Nv_j\operatorname{ReLU}
(u_j\cdot(1,x))$, and Gaussian energy
\[
 \kappa_N(\theta)=\frac dN\sum_j\|u_j\|^2+N\sum_jv_j^2.
\]
The prior has incoming coordinate variance $N\sigma^2/d$ and readout
variance $\sigma^2/N$. It is conditioned initially on
$|F(x)|>1/2$ at every cube input, and thereafter on exact labels.
For a Boolean function $g$, let $\kappa_N(g)$ be the infimum energy
of its strict global-margin realizations.

\begin{lemma}[A width-independent finite function count]
For $J\ge1$, put
\[
 V=\max\{1,(J+1)/(2\sqrt d)\},\quad
 T=\left\lceil128V^2D\log(4\,2^d)\right\rceil.
\]
Then the number of functions with $\kappa_N(g)\le J$ has logarithm at most
\[
 H(d,J)=\log(2+64DV)+T[\log2+D\log(3+128DV)].              \tag{2}
\]
In particular it has a polynomial majorant in $d,J$, independent of $N$.
\end{lemma}
\begin{proof}
Choose a realization of energy less than $J+1$. Its variation sum
$W=\sum_j|v_j|\|u_j\|$ is at most $(J+1)/(2\sqrt d)\le V$ by AM--GM.
Ignore zero atoms and write $F=W\E A$, where an atom is a signed ReLU
ridge with unit incoming norm. At every input $|A(x)|\le\sqrt D$.
For $T$ iid atoms, Hoeffding and a union bound on the $2^d$ inputs show
that with positive probability
\[
 \left\|\frac WT\sum_{j=1}^TA_j-F\right\|_\infty\le1/8.
\]
Indeed the failure probability is at most
$2\,2^d\exp[-T/(128V^2D)]\le1/2$.

Quantize every coordinate of the unit incoming vectors to a grid of
mesh $1/(64DV)$ within $[-1,1]$, including endpoints, with coordinate
error at most that mesh. Quantize $W\in[0,V]$ to mesh $1/(64D)$.
The incoming perturbation changes the output by at most $1/64$.
Quantized incoming coordinates remain in $[-1,1]$, so each quantized
feature has magnitude at most $D$, and amplitude quantization adds
at most $1/64$. The final error is less than $1/4$, preserving all
signs of the original strict-margin realization. The amplitude has
at most $2+64DV$ choices; each ordered atom has at most
$2(3+128DV)^D$ choices. Counting their descriptions proves (2).
\end{proof}

Suppose a target has a width-$C$ ReLU teacher, margin greater than
$1/2$, and all incoming and readout coordinates bounded by $B$.
Balancing each teacher neuron and padding zeros gives, at every $N\ge C$,
\[
 \kappa_N(f)\le J_T:=2\sqrt{dD}\,CB^2.
\]
For fixed finite $N$, the finite collection of realizable Boolean
functions satisfies the Gaussian small-noise formula
$2\sigma^2\log\mu_\sigma(U_{N,g})\to-\kappa_N(g)$.
This follows from an open ball about a near-minimizer for the lower
bound and the quadratic Gaussian tail for the upper bound. Therefore
for every $\beta>0$, all sufficiently small $\sigma>0$ satisfy
\[
 \pi(A^c)\le\beta\pi(f),\qquad
 A=\{g:\kappa_N(g)\le\kappa_N(f)\}.
\]
All ties are retained in $A$. Its log cardinality is bounded by
$H(d,\max\{1,J_T\})$. Equations (1)--(2) prove polynomial expected
lifetime mistakes for the actual strong-prior ReLU posterior under
two-sided comparison to uniform input, uniformly in the width budget.
The sufficient variance threshold may depend on each fixed $N$;
no width-uniform temperature threshold is asserted.

\section{Exponential mistakes with exponentially small upper-only shifts}
Now use the existing repository theorem
\emph{Superpolynomial online mistakes for isotropic Bayesian ReLU and softplus networks}
(the folder \texttt{isotropic-relu-allwidth}). For its parameter $k$, the
logical dimension is $m=6k+4$, the number of designated logical fibers is
$K=2^{k-1}$, and every physical fiber has size $F=2^{d-m}$.
Its target has a width-$\Theta(m^2)$ bounded-weight shallow ReLU teacher
and a size-$O(m)$ Boolean circuit. At each eligible fixed learner width,
the cold posterior prefers the wrong cheapest remaining star whenever
the entire calibration region and preceding fibers have been observed.
This statement concerns the original global absolute-margin support.

Let $C$ be the physical complement of all $K$ designated fibers.
Its size is $(2^m-K)F$. First draw uniformly from $C$ until every point
there has appeared. Then, in increasing star-cost order, draw uniformly
from $C\cup F_i$ until every point of fiber $F_i$ has appeared. Continue
with the next fiber; after the last phase retain its input law forever.
Repeated observations do not change a noiseless posterior.

\begin{theorem}[Tiny total variation does not replace a density lower bound]
For the above neural family there is a predictable input process with
$K=2^{k-1}$ designated first-arrival rounds for which, at every eligible
fixed width and all sufficiently small positive variance, posterior
majority is wrong on every designated round. Fresh posterior sampling
is wrong on all of them with probability at least $1-\eta$, for any
prescribed $\eta>0$ after choosing the variance threshold.

Writing $\mu$ for uniform physical input and $p=K/2^m=2^{-5k-5}$,
every conditional input law and every realized path satisfy
\[
 P_t\le\frac1{1-p}\mu,\qquad
 \TV(P_t,\mu)\le p,\qquad
 \sum_t\TV(P_{t+1},P_t)\le\frac p{1-p}.                  \tag{3}
\]
There is a deterministic horizon with $\log T=O(d)$ by which the
expected number of mistakes is $\Omega(K)$. Calibration and the horizon
are still enormous. The input laws have zeros, so they violate every
strictly positive density lower bound.
\end{theorem}
\begin{proof}
Every calibration and fiber-covering phase finishes almost surely.
Just before the first draw from $F_i$, the observed set is precisely
$C$ and all preceding whole fibers. Every minimum-energy completion
therefore mislabels every point of $F_i$, by the existing star theorem.
There are only finitely many stages and possible first points. Its
small-noise argument thus supplies one variance threshold for which
correct posterior probability is below $\eta/K$ at every such query.
Choose it also below $1/2$. This proves the majority statement and,
by conditional union bounds, the fresh-sampling statement.

The calibration law has density $2^m/(2^m-K)$ on $C$ and TV distance
$p$ from uniform. Every subsequent law is uniform on a set of size
$(2^m-K+1)F$, so its density upper bound is smaller and its TV distance
is $(K-1)/2^m$. The calibration-to-first-fiber change has TV
$1/(2^m-K+1)$, as does each subsequent change between disjoint new
fibers. There are $K$ changes. Thus their sum is
$K/(2^m-K+1)\le p/(1-p)$, proving (3) pathwise.

The mean completion time of the final phase is
\[
 T_0=|C|H_{|C|}+K(|C|+F)H_F,
\]
by the coupon-collector identity, including the replay delays from $C$.
Markov's inequality gives completion by $\lceil2T_0\rceil$ with
probability at least $1/2$. Taking, for example, $\eta=1/10$, the sampled
mistakes by that time have expectation at least $2K/5$; majority has
expectation at least $K/2$. Since $\log T_0=O(d)$, the asserted horizon
bound follows.
\end{proof}

\section{Interpretation and limits}
Both the positive shallow-ReLU result and the negative construction use
the same kind of independent Gaussian prior, exact label conditioning,
and global absolute-margin support. The positive density lower bound
forces continuing opportunities to eliminate every surviving wrong
function. Small TV distance and a near-unit upper density bound alone
can suppress a tiny set of those opportunities while exposing them in
an adversarial order. Their scarcity makes the instantaneous shift tiny,
without limiting the lifetime number of distinct mistakes.

The upper bound is a direct extension of existing finite-class and
neural-posterior arguments. The lower bound reuses the repository's
existing calibrated star geometry and its huge constants, rather than
providing a more practical example. The important open problem of
superpolynomial ordinary-ReLU mistakes with only observed margins is
untouched by this transfer.
\end{document}
